\section{连接不仅传位，还必须约定何时有效}

地址线、数据线和方向线即使全部接通，发起者仍要知道目标何时看到了请求，目标也要知道
处理器何时接收了响应。若启动存储一周期返回而 RAM 三周期返回，固定等待一个周期会让
处理器过早采样 RAM。教学互连因此使用请求/响应握手。

\subsection{请求通道的四个时序事实}

处理器把请求字段放稳后置 \code{req\_valid=1}。互连只有在能够接收时置
\code{req\_ready=1}。某个时钟边沿同时看到二者为 1，才表示请求已被接受。在此之前，
处理器必须保持地址、方向、宽度和写数据不变。

\begin{longtable}{|L{0.22\linewidth}|L{0.26\linewidth}|L{0.38\linewidth}|}
\hline
\rowcolor{BookTableHead}
信号 & 驱动者 & 为 1 时的含义 \\ \hline
\code{req\_valid} & 处理器 & 请求字段当前有效 \\ \hline
\code{req\_ready} & 互连 & 本边沿可以接受这份请求 \\ \hline
\code{resp\_valid} & 互连 & 响应状态和读数据当前有效 \\ \hline
\code{resp\_ready} & 处理器 & 本边沿可以消费响应 \\ \hline
\end{longtable}

响应通道采用同样规则。\code{resp\_valid} 与 \code{resp\_ready} 在同一边沿为 1，响应
才被消费。目标可以在请求接受后的任意较晚周期给出响应，处理器无需预先知道固定延迟。

\begin{center}
\begin{tikzpicture}[x=1.28cm,y=0.64cm]
\foreach \x in {0,1,2,3,4,5,6} {
  \draw[dashed,BookRule] (\x,0) -- (\x,5.7);
  \node[font=\scriptsize,BookMuted] at (\x,-0.25) {\x};
}
\foreach \y/\name in {5/req\_valid,4/req\_ready,3/resp\_valid,2/resp\_ready,1/state} {
  \node[anchor=east,font=\scriptsize] at (-0.25,\y) {\code{\name}};
}
\draw[very thick,BookAccent] (0,4.8)--(1,4.8)--(1,5.3)--(3,5.3)--(3,4.8)--(6,4.8);
\draw[very thick,BookAccent] (0,3.8)--(2,3.8)--(2,4.3)--(3,4.3)--(3,3.8)--(6,3.8);
\draw[very thick,BookAccent] (0,2.8)--(4,2.8)--(4,3.3)--(5,3.3)--(5,2.8)--(6,2.8);
\draw[very thick,BookAccent] (0,1.8)--(4,1.8)--(4,2.3)--(5,2.3)--(5,1.8)--(6,1.8);
\node[font=\scriptsize] at (2.5,0.95) {请求在边沿 2 接受};
\node[font=\scriptsize] at (4.5,0.55) {响应在边沿 4 消费};
\draw[flow] (2,1.25)--(2,3.55);
\draw[flow] (4,0.85)--(4,1.55);
\end{tikzpicture}
\end{center}
\diagramnote{处理器从周期 1 起保持请求，直到边沿 2 完成握手；目标稍后给出响应，并在边沿 4
被消费。有效信号单独为 1 都不表示传送已经发生。}

\subsection{为什么有效和就绪不能合成一根线}

若只有一根“完成”线，无法区分“发起者现在提供有效请求”和“接收者现在有容量接受”。
当互连忙于前一笔事务时，处理器必须保留请求；当处理器暂时不能接收响应时，互连也必须
保留响应。两端各自表达有效与就绪，才能在不同速度部件之间建立无歧义边界。

当前处理器每次只允许一笔未完成事务，因此不需要请求编号。若以后允许多笔并发且响应可能
乱序，就必须增加事务标识，并保存每笔请求对应的目标寄存器和 PC。那是性能优化引出的新
状态，不属于最小机器。

\subsection{写请求中的数据何时可以撤掉}

处理器从置 \code{req\_valid} 开始，就要同时保持 \code{write\_data} 和
\code{byte\_enable}。只有请求握手边沿之后，才可撤掉这些字段。请求被互连接受并不等于
目标写入已经完成；处理器还要等待响应握手，才能提交存储指令的顺序后继。

这形成两个边界：

\begin{enumerate}
\item 请求接受：互连已取得一份稳定请求，处理器可以释放请求通道；
\item 操作完成：目标已经执行或拒绝操作，响应被处理器消费。
\end{enumerate}

在简单零等待器件上，两者可能落在相邻甚至同一周期，但软件模型仍不应把它们定义成同一个
概念。

\subsection{物理连接中的方向与扇出}

教学图用箭头表达信息方向。地址、操作类型和写数据从处理器流向互连；读数据和响应状态从
目标流回处理器；目标选择信号从互连流向某个器件。时钟可以扇出到多个时序部件，复位也可
送达需要确定初态的控制状态。

\begin{center}
\begin{tikzpicture}[node distance=10mm and 17mm]
\node[stage,minimum width=3.4cm] (cpu) {处理器端口};
\node[stage,minimum width=3.5cm,right=of cpu] (fabric) {互连与译码};
\node[stage,minimum width=3.2cm,right=of fabric] (target) {被选目标};
\draw[flow] ([yshift=5mm]cpu.east)--node[above,font=\scriptsize]{地址、命令、写数据}([yshift=5mm]fabric.west);
\draw[flow] ([yshift=5mm]fabric.east)--node[above,font=\scriptsize]{目标请求}([yshift=5mm]target.west);
\draw[flow] ([yshift=-5mm]target.west)--node[below,font=\scriptsize]{目标响应}([yshift=-5mm]fabric.east);
\draw[flow] ([yshift=-5mm]fabric.west)--node[below,font=\scriptsize]{状态、读数据}([yshift=-5mm]cpu.east);
\end{tikzpicture}
\end{center}
\diagramnote{双向交互由两组单向信息构成。把一条线画成无说明的双向箭头，会隐藏谁必须保持
字段以及哪个边沿完成传送。}

\subsection{为什么地址不能由所有目标共同解释}

目标器件只需识别自己的内部偏移，不必都实现完整 32 位范围比较。互连先完成全局选择，再
把低位偏移交给目标。这减少重复逻辑，也把地址冲突集中在一处检查。

例如 RAM 收到内部偏移 \code{0x00002004}，无需知道全局地址曾是
\code{0x00102004}。但软件错误报告应保留原始全局地址，因为程序使用的是物理地址表，
不是器件内部偏移。

\subsection{复位连接不应清空所有 RAM}

处理器复位要确定 PC 和控制状态，串行控制器复位要清除 ready/overrun 等协议状态，互连
复位要清除未完成事务。RAM 内容是否在复位时清零是另一件事。多数 RAM 不会因处理器复位
自动擦除，固件必须把依赖初值的区域显式初始化。

这解释了为什么装入器不能在复位后看到旧 RAM 字节就判断任务仍有效。复位只重建控制起点，
不证明旧数据来自一次完整、校验通过的传输。

\subsection{三种更简单的启动方案为什么没有被选用}

\subsection{把任务直接固化在启动存储}

若机器永远只运行同一任务，可以让复位 PC 直接进入该任务，甚至省去串行控制器和装入器。
但任务更新需要重新写非易失存储，输入参数和可写栈仍需 RAM。本书需要观察“一个任务如何
被选择并获得初始现场”，所以保留固定固件与可替换任务的分离。

\subsection{从串行接口取一条执行一条}

处理器也可以等待四个指令字节到达，立即译码执行，不把完整映像存入 RAM。这样分支回到
先前指令时，字节已经从单向流中消失；函数调用和循环都要求随机访问旧指令。除非外部端
实现可寻址指令服务器和复杂协议，否则流式执行不能支持本章任务。

\subsection{边接收边从已到达区域执行}

固件可在前几条指令到达后就跳转，后续指令继续写入 RAM。任务若执行得比传输快，会取到
尚未到达的旧字节；若最终校验失败，部分代码已经产生外部效果。先接收、后验证、再一次性
提交 PC，牺牲了一点启动延迟，却给出清楚的执行边界。

\begin{longtable}{|L{0.23\linewidth}|L{0.28\linewidth}|L{0.35\linewidth}|}
\hline
\rowcolor{BookTableHead}
方案 & 减少的结构 & 失去的能力或边界 \\ \hline
任务固化在 ROM & 串行传输和装入器 & 任务不能方便替换 \\ \hline
串行流式执行 & 完整 RAM 映像 & 无法随机分支和重读指令 \\ \hline
接收与执行重叠 & 等待完整映像的时间 & 无法在执行前完成整体校验 \\ \hline
本章方案 & 无 & 需要 RAM、装入记录和一次明确移交 \\ \hline
\end{longtable}

比较的目的不是证明本章方案在所有机器上最好，而是指出它服务的目标：允许替换任务、允许
任意控制流，并且在执行前形成完整可验证映像。
